\scanpage{351}
i.e.
\[
(3)\qquad 2g'-2+\mathrm{card}\,I'=n(2g-2+\mathrm{card}\,I)
\]

(NB. \deleted{\unsure{\(n\)}} \((2g-2+\mathrm{card}\,I)\) est
\inserted{l'opposé de} la \unsure{car.} d'EP à supports
compacts de la courbe dont \(\unsure{\Pi|\pi}\) est le groupe
fondamental\ldots)

\margin{\uline{NB} dire qu'on \unsure{est} dans le cas
\unsure{anabélien} signifie que \(-\mathrm{EP}_{!}\) est \(\geq 1\)
--- cette \unsure{condition} est \unsure{donc conservée}
par passage \unsure{aux revêtements}.
L'entier \(-\mathrm{EP}_{!}\) \unsure{mesure} \unsure{positivité stricte}
\illegible{1} degré d'\unsure{anabélianité}, \unsure{en qq sorte}\ldots}

De façon évidente, \unsure{tte} discrétification de \(\unsure{\Pi|\pi}\)
en fournit une de \(\unsure{\Pi|\pi}'\) --- et il en est de
\(\overline{\mathrm m}\) pour les discrétifications orientées, modulo
un peu de cohomologie [qui revient à dire,
\inserted{dans le cas \unsure{discret} \deleted{\illegible{1}}}
que
\diagram{\begin{tikzcd}[ampersand replacement=\&]
H^2_{!}(\unsure{U|\mathcal U},\unsure{\mathbb Z|\mathbf Z}) \arrow[r] \arrow[d,phantom,sloped,"\simeq"] \& H^2_{!}(\unsure{U'|\mathcal U'},\unsure{\mathbb Z|\mathbf Z}) \arrow[d,phantom,sloped,"\simeq"] \\
T \& T'
\end{tikzcd}}
peut s'écrire sous la forme \(n\theta^{-1}\), où \(\theta\) est un iso. bien
déterminé (l'isom. trace) \unsure{\(\theta\)} \(T'\xrightarrow{\sim}T\).

Il n'est pas clair pour moi à première
vue si l'application
\[
\mathrm{Discret}(\unsure{\Pi|\pi})\longrightarrow\mathrm{Discret}(\unsure{\Pi|\pi}')
\]
\inserted{(pour \unsure{tt} \(\unsure{\Pi|\pi}'\) d'indice fini de \(\unsure{\Pi|\pi}\))}
est surjective, ou injective. L'injectivité
signifierait que : deux discrétifications \(\unsure{\Pi|\pi}_0,\unsure{\Pi|\pi}'_0\) de
\(\unsure{\Pi|\pi}\) qui sont commensurables, i.e. telles que
\(\unsure{\Pi|\pi}_0\cap\unsure{\Pi|\pi}'_0\) \unsure{soit} d'indice fini dans \(\unsure{\Pi|\pi}_0\) et dans \(\unsure{\Pi|\pi}'_0\),
\deleted{en fait, il suffit \unsure{qu'il le soit} dans l'un d'eux}
\inserted{égales}

\scanpage{352}
\mspage{178}
Supposons que \(\unsure{\Pi|\pi}'\) soit sous-groupe invariant
dans \(\unsure{\Pi|\pi}\), posons \(G=\unsure{\Pi|\pi}'/\unsure{\Pi|\pi}\)

\uline{Conjecture} L'application \(\mathrm{Discret}(\unsure{\Pi|\pi})\longrightarrow\mathrm{Discret}(\unsure{\Pi|\pi}')\)
est injective, et son image est formée des
discrétifications \(\unsure{\Pi|\pi}'_0\subset\unsure{\Pi|\pi}'\) telles que
\(G\longrightarrow\mathrm{Aut\,ext}_{\mathrm{lac}}(\unsure{\Pi|\pi}')
\simeq\mathrm{Aut\,ext}_{\mathrm{lac}}(\widehat{\unsure{\Pi|\pi}}'_0)
\simeq\widehat{\widehat{\mathcal T}}(\widehat{\unsure{\Pi|\pi}}'_0)\)
se factorise par \(\mathcal T(\unsure{\Pi|\pi}'_0)=\mathrm{Autext}_{\mathrm{lac}}(\unsure{\Pi|\pi}'_0)\)
(et en fait, \(\overline{\mathrm m}\) par \(\mathcal T(\unsure{\Pi|\pi}'_0)^+\)).

\margin{En d'autres termes, \(\forall g\in\unsure{\Pi|\pi}\),
\(\exists h\in\unsure{\Pi|\pi}'\) tel que \(\mathrm{int}(hg)(\unsure{\Pi|\pi}'_0)\subset\unsure{\Pi|\pi}'_0\)
i.e. \(\unsure{\unsure{\Pi|\pi}'|\unsure{\Pi|\pi}'_0}\cdot\unsure{\Pi|\pi}_0=\unsure{\Pi|\pi}\), où \(\unsure{\Pi|\pi}_0\) est le
normalisateur de \(\unsure{\Pi|\pi}'_0\) dans \(\unsure{\Pi|\pi}\).}

La condition pour qu'un \deleted{\illegible{1}} \(\unsure{\Pi|\pi}'_0\) soit dans
l'image est évidemment nécessaire. Montrons
qu'elle est suffisante, et que le \(\unsure{\Pi|\pi}_0^{\unsure{*}}\) donnant
\unsure{naissance} à \(\unsure{\Pi|\pi}'_0\) par \(\unsure{\Pi|\pi}_0\cap\unsure{\Pi|\pi}'\) est unique.
On a une action extérieure de \(G\) sur \(\unsure{\Pi|\pi}'_0\),
qui définit l'action extérieure sur \(\widehat{\unsure{\Pi|\pi}}'_0\simeq\unsure{\Pi|\pi}'\),
or \(\unsure{\Pi|\pi}\) se récupère canoniquement à partir de
cette dernière (car \(\mathrm{Centre}(\widehat{\unsure{\Pi|\pi}}')=1\)), et on
voit que c'est \unsure{le|l'} \deleted{\unsure{adhérence|extension} \unsure{dans|de}} \inserted{complété} profini de
l'extension \inserted{\(\unsure{\Pi|\pi}_0\)} de \(G\) par \(\unsure{\Pi|\pi}'_0\), définie par l'action
\inserted{ext.} de \(G\) sur \(\unsure{\Pi|\pi}'_0\). On a donc bien une
discrétification généralisée \inserted{\(\unsure{\Pi|\pi}_0\) de \(\unsure{\Pi|\pi}\)}, au sens de la
seule structure de \deleted{\unsure{gr}} groupe profini, et elle
induit \(\unsure{\Pi|\pi}'_0\) --- mais il faut voir qu'elle
est compatible avec la structure à lacets de \(\unsure{\Pi|\pi}\).

\scanpage{353}
Mais celle-ci s'explicite, à partir de la
structure à lacets de \(\unsure{\Pi|\pi}'\), on prend les
sous-groupes à lacets \inserted{\(L'\)} de \(\unsure{\Pi|\pi}'\), et leurs
Centralisateurs dans \(\unsure{\Pi|\pi}\). Faisons \unsure{itou} pour
\(\unsure{\Pi|\pi}_0\unsure{\supset|>}\unsure{\Pi|\pi}'_0\), on devrait trouver une structure
à lacets sur \(\unsure{\Pi|\pi}_0\), donnant lieu aux
\(\overline{\mathrm m}\ \unsure{d_i|\delta_i}\).
J'ai l'impression que ça ne
doit pas être très vache à prouver --- ce
serait \unsure{connu} si on savait déjà que \unsure{tte}
extension \inserted{``sans torsion''} d'un groupe fini par un
groupe à lacets est (de la façon canonique)
un groupe à lacets\ldots

\uline{Corollaire} En tous cas, l'application
\(\mathrm{Discret}(\unsure{\Pi|\pi})\longrightarrow\mathrm{Discret}(\unsure{\Pi|\pi}')\)
est \uline{injective}.
\figure{[Trait vertical à gauche de la ligne commençant par \(\mathrm{Discret}(\unsure{\Pi|\pi})\).]}

\uline{NB} Elle ne doit pas tjs être bijective,
car il doit y avoir \unsure{une|des} \unsure{action|actions} \inserted{\unsure{extérieure|extérieures}} d'un
\(G\) dans un \(\widehat{\unsure{\Pi|\pi}}'_0\) i.e. des hom
\(G\longrightarrow\widehat{\widehat{\mathcal T}}(\widehat{\unsure{\Pi|\pi}}'_0)\), qui ne se \unsure{factorise|factorisent} pas
par \(\mathcal T(\unsure{\Pi|\pi}'_0)\), et qui pourtant est aussi
\unsure{bonne|bonnes} du pt de vue groupe profini que
\unsure{celle|celles} provenant d'un \(\unsure{\Pi|\pi}\) et d'un \unsure{ss-groupe}

\scanpage{354}
\mspage{179}

invariant $\unsure{\Pi|\pi}'$. En fait, partons d'une telle situation \uline{discrète} $\unsure{\Pi|\pi}_0 \unsure{\supset|>} \unsure{\Pi|\pi}'_0$, d'où $G=\unsure{\Pi|\pi}_0/\unsure{\Pi|\pi}'_0$ et $G\longrightarrow \operatorname{Autext}_{\mathrm{lac}}(\unsure{\Pi|\pi}'_0)\deleted{\illegible{1}}=\mathcal T(\unsure{\Pi|\pi}'_0)\subset\unsure{\widehat{\mathcal T}|\widehat{\widehat{\mathcal T}}}(\widehat{\unsure{\Pi|\pi}}'_0)$, et \uline{conjuguons} $G$ par un élément \inserted{$\gamma$} de $\widehat{\widehat{\mathcal T}}(\unsure{\Pi|\pi}'_0)$, \unsure{\uline{de façon}|\inserted{de façon}\deleted{\illegible{2}}} à trouver $G\longrightarrow\widehat{\mathcal T}(\widehat{\unsure{\Pi|\pi}}'_0)$ qui ne se \unsure{factorise} pas par $\mathcal T(\unsure{\Pi|\pi}^{\unsure{'}}_0)$. On trouve une extension \inserted{$\unsure{\widehat{\unsure{\Pi|\pi}}|\unsure{\Pi|\pi}}^{\unsure{\natural|\sharp}}$} de $G$ par $\widehat{\unsure{\Pi|\pi}}'_0$, \uline{isomorphe} à $\widehat{\unsure{\Pi|\pi}}_0$ (en tant qu'extension de $G$ par $\widehat{\unsure{\Pi|\pi}}'_0$), donc la structure à lacets de $\widehat{\unsure{\Pi|\pi}}'_0$ en définit une sur $\unsure{\Pi|\pi}^{\unsure{\natural|\sharp}}$ \deleted{et est} \inserted{(de façon à être} induite par cette dernière), mais pourtant la discrétification \uline{choisie} \inserted{$\unsure{\Pi|\pi}_0$} de $\widehat{\unsure{\Pi|\pi}}_0$ n'est pas induite par une de $\unsure{\Pi|\pi}$. Pour

$\ast$ Faire la construction, il suffit de trouver un élément \inserted{$\gamma\in\widehat{\widehat{\mathcal T}}(\unsure{\Pi|\pi}'_0)$ \deleted{\unsure{fini}} tel que $\operatorname{int}(\gamma).G\unsure{\not\subset|\notin}\mathcal T$,} \deleted{\unsure{$u$} \unsure{d'ordre} \illegible{1} \unsure{provenant} pas de, \unsure{dans} $\mathcal T(\unsure{\Pi|\pi}'_0)$, \unsure{ou se} \unsure{relevant} par un élément d'ordre fini de $\unsure{\mathcal G|\mathcal C}(\unsure{\Pi|\pi}'_0)^+=\operatorname{Aut}_{\mathrm{lac}}(\unsure{\Pi|\pi}'_0)^+$, et \deleted{dont la \unsure{conjugaison} dans $\widehat{\widehat{\mathcal T}}$ \unsure{un}} \inserted{un $\gamma\in\widehat{\widehat{\mathcal T}}(\unsure{\Pi|\pi}'_0)$, tel que} $\gamma u\gamma^{-1}\notin\mathcal T(\unsure{\Pi|\pi}'_0)$, \inserted{\deleted{i.e. tel que $u$ ne \unsure{fixe}}} \inserted{$\gamma^{-1}.\unsure{\Pi|\pi}'_0\subset\widehat{\unsure{\Pi|\pi}}'_0$.}}
\inserted{i.e. tel que $G$ \unsure{ne stabilise pas} $\gamma^{-1}(\unsure{\Pi|\pi}'_0)\subset\widehat{\unsure{\Pi|\pi}}'_0$}

\deleted{\illegible{1}} Un tel $\gamma$ existe \unsure{tjrs}....

\scanpage{355}

Bien sûr, il n'y a pas de raison, \deleted{que} pour \unsure{des|deux} discrétifications $\unsure{\Pi|\pi}'_0,\unsure{\Pi|\pi}_1$ de $\unsure{\Pi|\pi}$, que si $\unsure{\Pi|\pi}_0,\unsure{\Pi|\pi}_1$ sont $\unsure{\Pi|\pi}$-conjugués (i.e. définissent \deleted{des} \inserted{la $\overline{\mathrm m}$} \unsure{prédiscrt.|prédiscrét.} \unsure{strictes|stricte} de $\unsure{\Pi|\pi}$), il en soit de $\overline{\mathrm m}$ pour $\unsure{\Pi|\pi}'_0$ et $\unsure{\Pi|\pi}'_1$ pour l'action intérieure de $\unsure{\Pi|\pi}'$. (Il faudrait que $\unsure{\Pi|\pi}_0,\unsure{\Pi|\pi}_1$ soient conjugués par $\unsure{\Pi|\pi}'$, et pas seulement par $\unsure{\Pi|\pi}$). Par contre, je présume que si $\unsure{\Pi|\pi}_0,\unsure{\Pi|\pi}_1$ sont ``adhérents'' l'un à l'autre dans l'espace des \unsure{discrétifications|discrétific.} de $\unsure{\Pi|\pi}$, i.e. s'ils définissent une $\overline{\mathrm m}$ prédiscrétification, il en va de $\overline{\mathrm m}$ pour $\unsure{\Pi|\pi}'_0,\unsure{\Pi|\pi}'_1$, et que l'application
\[
\tag{4}
\operatorname{Prédiscrt}^{\unsure{+|\ast}}(\unsure{\Pi|\pi})\longrightarrow\operatorname{Prédiscrt}^{\unsure{+|\ast}}(\unsure{\Pi|\pi}')
\]
qu'on obtient ainsi, est encore bijective, et qu'elle passe à son tour au quotient, pour définir une application bijective
\[
\tag{5}
\unsure{\operatorname{Préarithm}}(\unsure{\Pi|\pi})\longrightarrow\unsure{\operatorname{Préarithm}}(\unsure{\Pi|\pi}'),
\]
et que si $\Sigma,\Sigma'$ se correspondent par cette dernière, qu'on a un isom canonique
\[
\tag{6}
\unsure{\mathbb{\Gamma}|\Gamma}_{\Sigma}\longrightarrow\unsure{\mathbb{\Gamma}|\Gamma}_{\Sigma'},
\]
\unsure{comp.$^{\mathrm t}$ble},

\scanpage{356}
\mspage{180}

avec les actions de ces groupes sur l'ens $P_{\Sigma}$ des prédiscrétifications de $\unsure{\Pi|\pi}$ \unsure{compatibles|comp.tibles} à $\Sigma$, et l'ens \inserted{$P_{\Sigma'}$} des prédiscrétifications de $\unsure{\Pi|\pi}'$ \unsure{compatibles|comp.tibles} à $\Sigma'$ (qui sont resp. des torseurs à g. sous $\unsure{\mathbb{\Gamma}|\Gamma}_{\Sigma},\unsure{\mathbb{\Gamma}|\Gamma}_{\Sigma'}$). En même temps, on trouvera \inserted{donc} que $\Sigma$ est \unsure{unique} pour $\unsure{\Pi|\pi}$ i.e. invariant dans $\widehat{\widehat{\mathcal T}}(\unsure{\Pi|\pi})$, ssi $\Sigma'$ est \unsure{unique} dans $\unsure{\Pi|\pi}'$ i.e. invariant dans $\widehat{\widehat{\mathcal T}}(\unsure{\Pi|\pi}')$.

En tous cas, la conjecture fondamentale qui \unsure{interprète} les $\unsure{\mathbb{\Gamma}|\Gamma}_{g,\unsure{\nu|v}}$ comme $\unsure{\mathbb{\Gamma}|\Gamma}_0=\operatorname{Gal}(\overline{\mathbb Q}_0/\mathbb Q)$ \unsure{aurait|aura} au moins comme conséquence que l'application $\operatorname{Discr}^{+}(\unsure{\Pi|\pi})\unsure{-|\longrightarrow}\operatorname{Discr}^{+}(\unsure{\Pi|\pi}')$ passe au quotient \inserted{de} deux façons, pour donner un diagramme comm.

\diagram{
\begin{tikzcd}[row sep=large,column sep=large,ampersand replacement=\&]
\operatorname{Discr}^{+}(\unsure{\Pi|\pi}) \arrow[r,no head] \arrow[d]
  \& \operatorname{Discr}^{+}(\unsure{\Pi|\pi}') \arrow[d] \\
\unsure{P|p}=\unsure{\operatorname{Prédiscr}}^{+}(\unsure{\Pi|\pi}) \arrow[r] \arrow[d]
  \& \unsure{P|p}'=\operatorname{Prédiscr}^{+}(\unsure{\Pi|\pi}') \arrow[d] \\
A=\operatorname{Arithm}(\unsure{\Pi|\pi}) \arrow[r]
  \& A'=\operatorname{Arithm}(\unsure{\Pi|\pi}')
\end{tikzcd}
}

et que le \deleted{\unsure{diagr|diagr.}} carré inférieur est cartésien (sans préjuger de l'\unsure{injectivité} ou de la \unsure{bijectivité} de l'application $A\longrightarrow A'$). \unsure{Donc} il faudrait

\scanpage{357}
que \deleted{\unsure{déjà}} si $\unsure{\Pi|\pi}_0,\unsure{\Pi|\pi}_1$ sont deux discrétifications \inserted{(\unsure{orientées})} de $\unsure{\Pi|\pi}$, qui définissent la $\overline{\mathrm m}$ prédiscrétification (\unsure{orientée}), alors il en va de $\overline{\mathrm m}$ pour \inserted{leurs \unsure{traces} sur $\unsure{\Pi|\pi}'$,} $\unsure{\Pi|\pi}'_0$ \unsure{et} $\unsure{\Pi|\pi}'_1$, relativement à $\unsure{\Pi|\pi}'$. Mais déjà si $\unsure{\Pi|\pi}_0$ et $\unsure{\Pi|\pi}_1$ sont $\unsure{\Pi|\pi}$-conjugués (i.e. définissent $\overline{\mathrm m}$ discrét. (\unsure{orientée}, \uline{\unsure{stricte}})), ce n'est pas tellement clair --- \deleted{\illegible{3}} mais si, car on aura $\unsure{\Pi|\pi}=\unsure{\Pi|\pi}'\cdot\unsure{\Pi|\pi}_0$ donc si $\unsure{\Pi|\pi}_1=\operatorname{int}(u)\unsure{\Pi|\pi}_0$ écrivant $u=u'u_0$ avec $u'\in\unsure{\Pi|\pi}'$, $u_0\in\unsure{\Pi|\pi}_0$, on aura
\[
\unsure{\Pi|\pi}_1=\operatorname{int}(u')\operatorname{int}(u_0)\unsure{\Pi|\pi}_0=\operatorname{int}(u')\unsure{\Pi|\pi}_0,
\]
donc $\unsure{\Pi|\pi}'_1=\operatorname{int}(u')\unsure{\Pi|\pi}'_0$ donc $\unsure{\Pi|\pi}'_0$ et $\unsure{\Pi|\pi}'_1$ sont conjugués.

Tâchons de procéder de $\overline{\mathrm m}$ dans le cas général d'un $u\in\widehat{\unsure{\mathcal{C}|\mathcal{G}}(\unsure{\Pi|\pi}_0)}^+$, \inserted{tel que $\unsure{\Pi|\pi}_1=u(\unsure{\Pi|\pi}_0)$,} en écrivant si possible $u=u'u_0$, avec $u'\in\widehat{\unsure{\mathcal{C}|\mathcal{G}}(\unsure{\Pi|\pi}_0,\unsure{\Pi|\pi}'_0)}^+$, et $u_0\in\unsure{\mathcal{C}|\mathcal{G}}(\unsure{\Pi|\pi}_0)^+$, où on définit $\unsure{\mathcal{C}|\mathcal{G}}(\unsure{\Pi|\pi}_0,\unsure{\Pi|\pi}'_0)^+$ comme le groupe des \unsure{i-automs|l-automs} \inserted{\unsure{directs|discrets}} de $\unsure{\Pi|\pi}_0$ qui \uline{stabilisent $\unsure{\Pi|\pi}'_0$}, et $\widehat{\unsure{\mathcal{C}|\mathcal{G}}}(\unsure{\Pi|\pi}_0,\unsure{\Pi|\pi}'_0)^+$ est son \unsure{compactifié} profini. On aurait alors $\unsure{\Pi|\pi}_1=u'(u_0(\unsure{\Pi|\pi}_0))=\dot{u}'(\unsure{\Pi|\pi}_0)$, où \deleted{\illegible{1}} $\dot{u}'$ désigne l'image de $u'$ dans $\widehat{\unsure{\mathcal{T}|\mathcal{C}}(\unsure{\Pi|\pi}'_0)}$. En tous cas, il n'y a aucun Pb si $\unsure{\Pi|\pi}'$ est un ss-groupe \deleted{\unsure{fermé}} \uline{caractéristique} de $\unsure{\Pi|\pi}$, car alors on a un hom

\scanpage{358}
\mspage{181}
discret $\unsure{\mathcal{C}|\mathcal{G}}(\unsure{\Pi|\pi}_0)\longrightarrow\unsure{\mathcal{C}|\mathcal{G}}(\unsure{\Pi|\pi}'_0)$, définissant \unsure{l'|le} $\unsure{\widehat{\unsure{\mathcal{C}|\mathcal{G}}}(\widehat{\unsure{\Pi|\pi}}_0)}\longrightarrow\unsure{\widehat{\unsure{\mathcal{C}|\mathcal{G}}}(\widehat{\unsure{\Pi|\pi}}'_0)}$. Comme les sous groupes \unsure{ouverts} caract. de $\unsure{\Pi|\pi}$ sont cofinaux, on est ramené (pour les questions de factorisabilité \inserted{de $D(\unsure{\Pi|\pi})\longrightarrow D(\unsure{\Pi|\pi}')$ en $P(\unsure{\Pi|\pi})\longrightarrow P(\unsure{\Pi|\pi}')$ et $A(\unsure{\Pi|\pi})\longrightarrow A(\unsure{\Pi|\pi}')$} et de bijectivité des applications $\unsure{P}(\unsure{\Pi|\pi})\longrightarrow P(\unsure{\Pi|\pi}')$ et $A(\unsure{\Pi|\pi})\longrightarrow A(\unsure{\Pi|\pi}')$) on est ramené au cas où $\unsure{\Pi|\pi}'$ est ss-groupe \uline{caractéristique}. On a \unsure{bien|vu} \inserted{\unsure{alors}} la factorisabilité de $D(\unsure{\Pi|\pi})\longrightarrow D(\unsure{\Pi|\pi}')$ en $P(\unsure{\Pi|\pi})\longrightarrow P(\unsure{\Pi|\pi}')$, \inserted{\unsure{nous voulons}} \unsure{montrer|montrons} que cette application est injective \unsure{...} \inserted{Cela} \deleted{\unsure{cela signifie}}
\inserted{\textcircled{a}} signifie que l'image inverse dans $\widehat{\widehat{\unsure{\mathcal{C}|\mathcal{G}}}}(\widehat{\unsure{\Pi|\pi}}_0)$ de $\widehat{\unsure{\mathcal{C}|\mathcal{G}}}(\widehat{\unsure{\Pi|\pi}}'_0)^+$ est $\widehat{\unsure{\mathcal{C}|\mathcal{G}}}(\widehat{\unsure{\Pi|\pi}}_0)^+$, la surjectivité \inserted{\textcircled{b}} signifie que tt élément de $\widehat{\widehat{\unsure{\mathcal{C}|\mathcal{G}}}}(\widehat{\unsure{\Pi|\pi}}'_0)$ est congru mod $\widehat{\unsure{\mathcal{C}|\mathcal{G}}(\unsure{\Pi|\pi}'_0)}^+$ à un élément de l'image de $\widehat{\widehat{\unsure{\mathcal{C}|\mathcal{G}}}}(\unsure{\Pi|\pi}_0)$. La factorisabilité en $A(\unsure{\Pi|\pi})\longrightarrow A(\unsure{\Pi|\pi}')$ signifie que par \inserted{\textcircled{c}} l'application précédente $\widehat{\widehat{\unsure{\mathcal{C}|\mathcal{G}}}}(\unsure{\Pi|\pi}_0)\hookrightarrow\widehat{\widehat{\unsure{\mathcal{C}|\mathcal{G}}}}(\unsure{\Pi|\pi}'_0)$, (NB il est immédiat que \unsure{c'est} injectif) \deleted{\unsure{est}} envoie
\[
M(\unsure{\Pi|\pi}_0)=\operatorname{Norm}_{\unsure{\widehat{\widehat{\mathcal{C}}}|\widehat{\mathcal{C}}}(\unsure{\Pi|\pi}_0)}\bigl(\widehat{\unsure{\mathcal{C}|\mathcal{G}}}^{+}(\unsure{\Pi|\pi}_0)\bigr)
\]
dans $M(\unsure{\Pi|\pi}'_0)$, l'injectivité de $A(\unsure{\Pi|\pi})\longrightarrow A(\unsure{\Pi|\pi}')$ que

\scanpage{359}

\textcircled{d} l'on a même \unsure{que} $M(\unsure{\Pi|\pi})$ est
\unsure{exactement|ex.\textsuperscript{t}} l'image inverse de $M(\unsure{\Pi|\pi}')$,
la \unsure{surjectivité} de
$A(\unsure{\Pi|\pi})\unsure{\twoheadrightarrow|\longrightarrow}A(\unsure{\Pi|\pi}')$ que
\inserted{\textcircled{e}} tt élément de
$\widehat{\widehat{\unsure{\mathcal C|\mathcal G}}}(\unsure{\Pi|\pi}'_0)$
est congru mod $M(\unsure{\Pi|\pi}'_0)$ à un élément de
$\widehat{\widehat{\unsure{\mathcal C|\mathcal G}}}(\unsure{\Pi|\pi}_0)$
\margin{plus faible que (b)}
--- et ces conditions \unsure{réunies} impliquent que l'hom canonique
\diagram{\begin{tikzcd}[ampersand replacement=\&]
\mathbb{\Gamma}_{\deleted{\illegible{1}}\unsure{\Pi|\pi}_0}
  \arrow[r] \arrow[d, equal]
  \& \mathbb{\Gamma}_{\unsure{\Pi|\pi}'_0} \arrow[d, equal] \\
M(\unsure{\Pi|\pi}_0)/\widehat{\unsure{\mathcal C|\mathcal G}}(\unsure{\Pi|\pi}_0)
  \& M^{\unsure{\sharp|*}}(\unsure{\Pi|\pi}'_0)/\widehat{\unsure{\mathcal C|\mathcal G}}(\unsure{\Pi|\pi}'_0)
\end{tikzcd}}
est iso -- il suffit $\overline{\mathrm m}$ pour ceci d'avoir
\unsure{(c)|(e)} (pour pouvoir définir cette application) et (a)
(pour son injectivité), et (b) et le renforcement (d) de
\unsure{(c)|(e)} (pour sa bijectivité. Donc on voit que
\unsure{tout} est suspendu aux propriétés de l'hom d'inclusions de groupes
\diagram{[Diagramme (7). Deux rangées, portant respectivement
$\widehat{\unsure{\mathcal C|\mathcal G}}^{+}(\unsure{\Pi|\pi}_0)
\subset M(\unsure{\Pi|\pi}_0)\subset
\widehat{\widehat{\unsure{\mathcal C|\mathcal G}}}(\unsure{\Pi|\pi}_0)$
et
$\widehat{\unsure{\mathcal C|\mathcal G}}^{+}(\unsure{\Pi|\pi}'_0)
\subset M(\unsure{\Pi|\pi}'_0)\subset
\widehat{\widehat{\unsure{\mathcal C|\mathcal G}}}(\unsure{\Pi|\pi}'_0)$.
Entre les termes de gauche, et entre ceux de droite, traits verticaux
recourbés à leur extrémité supérieure; les pointes ne sont pas nettement
discernables. Entre les deux $M$, flèche vers le bas, recourbée en haut,
avec un point d'interrogation à droite; cette flèche et le point
d'interrogation sont entourés de crochets.]}
à charge de prouver a) b) (qui ne \unsure{concernent} que le carré composé)
et \unsure{c)|e)}.

\deleted{D'ailleurs d) \unsure{sera} conséquence
\inserted{\unsure{immédiate}} de a), et du fait
(qui \unsure{assurera|implique} \unsure{c)|e)}) que tt élément de
$\widehat{\unsure{\mathcal C|\mathcal G}}^{+}(\unsure{\Pi|\pi}_0)$ est congru mod
$\unsure{\mathcal C|\mathcal G}$}

\scanpage{360}
\mspage{182}

\unsure{Peut-être|Peut. être} les propriétés a) b) c) sont elles
tout à fait fausses -- pourtant a) (qui correspond à l'injectivité de
$P(\unsure{\Pi|\pi})\longrightarrow P(\unsure{\Pi|\pi}')$) semble \unsure{assez} plausible.
La propriété b) (qui exprimerait la surjectivité de
$P(\unsure{\Pi|\pi})\longrightarrow P(\unsure{\Pi|\pi}')$) est beaucoup plus problématique,
voire douteuse. L'ennui, c'est que le groupe
$\widehat{\widehat{\unsure{\mathcal C|\mathcal G}}}(\unsure{\Pi|\pi}_0)$
peut être ``beaucoup plus petit'' que
$\widehat{\widehat{\unsure{\mathcal C|\mathcal G}}}(\unsure{\Pi|\pi}'_0)$,
on s'en rend compte en \unsure{passant} au quotient par le sous-groupe
$\widehat{\unsure{\Pi|\pi}}'_0=\unsure{\Pi|\pi}'$ (contenu dans le plus petit des groupes de (7),
et invariant dans le plus grand), on trouve sur la deuxième ligne
les groupes $\widehat{\mathcal T}^{+}(\unsure{\Pi|\pi}'_0)$,
\deleted{\illegible{1}} $\unsure{\mathcal N|N}(\unsure{\Pi|\pi}'_0)$ et
$\widehat{\widehat{\mathcal T}}(\unsure{\Pi|\pi}'_0)$, sur la première
des extensions des groupes \unsure{corresp.dants|correspondants}
pour $\unsure{\Pi|\pi}_0$ ($\widehat{\mathcal T}^{+}(\unsure{\Pi|\pi}_0)$,
$\unsure{\mathcal N|N}(\unsure{\Pi|\pi}_0)$ et
$\widehat{\widehat{\mathcal T}}(\unsure{\Pi|\pi}_0)$) par
\inserted{le groupe fini} $G=\unsure{\Pi|\pi}_0/\unsure{\Pi|\pi}'_0\simeq\unsure{\Pi|\pi}/\unsure{\Pi|\pi}'$,
\deleted{\unsure{soient|soit}
$\unsure{\widetilde{\widehat{\mathcal T}}^{+}|\widetilde{\mathcal T}^{+}}$
ou $\mathcal T$}
\unsure{notées|notés} par des $\sim$, de sorte que \unsure{.|--}
\diagram{[Diagramme (8). Rangée supérieure:
$G\subset
\widetilde{\widehat{\unsure{\mathcal T|\mathcal C}}}^{+}(\unsure{\Pi|\pi}_0)
\subset\widetilde{\unsure{\mathcal N|N}}(\unsure{\Pi|\pi}_0)
\subset\widetilde{\widehat{\widehat{\mathcal T}}}(\unsure{\Pi|\pi}_0)$.
Le premier symbole après $G$ est surchargé, et deux petits points sont
visibles entre son $+$ et la parenthèse. Rangée inférieure:
$\widehat{\mathcal T}^{+}(\unsure{\Pi|\pi}'_0)
\subset\unsure{\mathcal N|N}(\unsure{\Pi|\pi}'_0)
\subset\widehat{\widehat{\mathcal T}}(\unsure{\Pi|\pi}'_0)$.
Flèche oblique, recourbée à son départ, de $G$ vers
$\widehat{\mathcal T}^{+}(\unsure{\Pi|\pi}'_0)$.
Entre les deux termes portant un $+$, trait vertical recourbé en haut,
sans pointe nettement discernable. À droite, flèche verticale vers le bas,
recourbée en haut. Entre les deux $\unsure{\mathcal N|N}$, flèche vers le bas,
recourbée en haut, avec un point d'interrogation à droite;
cette flèche et le point d'interrogation sont entourés de crochets.]},
qui montre que les groupes de la première ligne
